%Version 3.1 December 2024
% See section 11 of the User Manual for version history
%
%%%%%%%%%%%%%%%%%%%%%%%%%%%%%%%%%%%%%%%%%%%%%%%%%%%%%%%%%%%%%%%%%%%%%%
%%                                                                 %%
%% Please do not use \input{...} to include other tex files.       %%
%% Submit your LaTeX manuscript as one .tex document.              %%
%%                                                                 %%
%% All additional figures and files should be attached             %%
%% separately and not embedded in the \TeX\ document itself.       %%
%%                                                                 %%
%%%%%%%%%%%%%%%%%%%%%%%%%%%%%%%%%%%%%%%%%%%%%%%%%%%%%%%%%%%%%%%%%%%%%

%%\documentclass[referee,sn-basic]{sn-jnl}% referee option is meant for double line spacing

%%=======================================================%%
%% to print line numbers in the margin use lineno option %%
%%=======================================================%%

\documentclass[lineno,referee,pdflatex,sn-basic]{sn-jnl}% Basic Springer Nature Reference Style/Chemistry Reference Style

%%%% Standard Packages
%%<additional latex packages if required can be included here>

\usepackage{graphicx}%
\usepackage{multirow}%
\usepackage{amsmath,amssymb,amsfonts}%
\usepackage{amsthm}%
\usepackage{mathrsfs}%
\usepackage[title]{appendix}%
\usepackage{xcolor}%
\usepackage{textcomp}%
\usepackage{manyfoot}%
\usepackage{booktabs}%
\usepackage{algorithm}%
\usepackage{algorithmicx}%
\usepackage{algpseudocode}%
\usepackage{listings}%
%%%%

%%%%%=============================================================================%%%%
%%%%  Remarks: This template is provided to aid authors with the preparation
%%%%  of original research articles intended for submission to journals published 
%%%%  by Springer Nature. The guidance has been prepared in partnership with 
%%%%  production teams to conform to Springer Nature technical requirements. 
%%%%  Editorial and presentation requirements differ among journal portfolios and 
%%%%  research disciplines. You may find sections in this template are irrelevant 
%%%%  to your work and are empowered to omit any such section if allowed by the 
%%%%  journal you intend to submit to. The submission guidelines and policies 
%%%%  of the journal take precedence. A detailed User Manual is available in the 
%%%%  template package for technical guidance.
%%%%%=============================================================================%%%%



\begin{document}

\title[Article Title]{Haemocytes of Crassostrea virginica as a dual-purpose source of ultra-long DNA for genome scaffolding and environmental DNA surveillance}

%%=============================================================%%
%% GivenName	-> \fnm{Joergen W.}
%% Particle	-> \spfx{van der} -> surname prefix
%% FamilyName	-> \sur{Ploeg}
%% Suffix	-> \sfx{IV}
%% \author*[1,2]{\fnm{Joergen W.} \spfx{van der} \sur{Ploeg} 
%%  \sfx{IV}}\email{iauthor@gmail.com}
%%=============================================================%%

\author[1]{\fnm{Alyssa} \sur{Strickland}}

\author[2]{\fnm{Kshitij} \sur{Amar}}

\author[1]{\fnm{Bo-Young} \sur{Lee}}

\author[1]{\fnm{Caylin} \sur{Grove}}

\author[2]{\fnm{Miten} \sur{Jain}}

\author*[1]{\fnm{Bonnie L.} \sur{Brown}} \email{bonnie.brown@unh.edu}

\affil*[1]{\orgdiv{Department of Biological Sciences}, \orgname{University of New Hampshire}, \orgaddress{\city{Durham}, \postcode{03874}, \state{NH}, \country{USA}}}

\affil[2]{\orgdiv{Department of Bioengineering}, \orgname{Northeastern University}, \orgaddress{\city{Boston}, \postcode{02115}, \state{MA}, \country{USA}}}

%%==================================%%
%% Sample for unstructured abstract %%
%%==================================%%

\abstract{Sequencing and assembling ultra-long DNA (ULDNA) scaffolds can extend utility of the existing eastern oyster \textit{Crassostrea virginica}, reference genome. In an attempt to create a long-read oyster reference, methodology to successfully isolate ULDNA from eastern oysters was optimized to sequester ultralong oyster DNA from mucopolysaccharides while avoiding excessive shearing. Analysis using a PromethION\texttrademark{} flow cell yielded read lengths (average N50 = 16.5kb, longest read 508kb) that were comparable to what was expected. After quality filtering ($Q \geq 9$), the number of reads $393 \times 10^{3}$ were mapped to 91.4\% of the reference genome (649Mbp, N50=14kb). However, the sequencing coverage was not sufficiently deep to accomplish de novo assembly. Of all ULDNA reads, 94.6\% were successfully mapped to \textit{C. virginica}. Most of the unmapped reads (5.4\% of all reads) were classified to bacteria (56\%) and human (33\%), revealing, respectively, an unexpected oyster pathogen and signaling sewage contamination in the vicinity of oyster harvest.}

\maketitle

\section{Introduction}\label{sec1}
Eastern oysters, \textit{Crassostrea virginica}, are both environmentally and economically important in the estuaries where they reside. Eastern oysters provide reef structures that act as habitats for juvenile marine organisms and aid in protection from predators at vulnerable life stages \cite{coen_role_1999, harding_oyster_2001}. Additionally, oyster reef structures aid in shoreline protection by reducing storm event wave energy \cite{brandon_evidence_2016, la_peyre_assessing_2015, scyphers_oyster_2011}. 

Perhaps the most notable function of eastern oysters is as water filterers. At their pre-decline population levels in Chesapeake Bay, eastern oysters were most likely able to filter the water volume from their entire estuarine systems in a few days \cite{newell_ecological_1988}. Presently, the universal oyster filtration potential is significantly reduced due to their population decline, resulting in pronounced negative effects to the ecosystem \cite{zu_ermgassen_quantifying_2013}. Oysters filter to consume their food source, phytoplankton, and in the process, they remove other microbes, improve water clarity, and transfer biomass to the estuary floor where in the instance of oyster reefs, they promote more favorable nutrient cycling and ultimately decreasing levels of nitrogen pollution \cite{bricker_bioextractive_2020, dame_carbon_1989, ehrich_review_2015, grizzle_seston_2008, grizzle_growth_2017, higgins_nutrient_2011, higgins_effect_2013, hoellein_effect_2014, kellogg_use_2014, newell_influence_2005}. Since they are continuously filtering at high rates, it can be hypothesized that eastern oysters can act as “natural eDNA samplers.” Studies have documented the reduction in nitrogen levels due to oyster filtration (Bricker et al., 2018; Kellogg et al., 2014), bringing forward the idea that bivalve species, such as the eastern oyster, could be used as an ecosystem engineer to improve nutrient levels in the estuaries where they are found. However, with a dramatic decline in eastern oyster populations (Bersoza Hernández et al., 2018a; R. Grizzle & Ward, 2016; Mackenzie, 2007; Rothschild et al., 1994), there has been a substantial loss of the ecosystem services that they provide (Zu Ermgassen et al., 2013).
A better understanding of eastern oyster responses to their environment is resulting from advances of omics-techniques. Genomic studies have supported extensive population analyses (Hornick & Plough, 2022; A. Strickland & Brown, 2024; Varney et al., 2009) transcriptomic studies have uncovered expression patterns in oyster as they respond to changing environmental conditions, stress, and disease (Chapman et al., 2011; Eierman & Hare, 2014, 2016; S. Wang et al., 2010). Sequencing and annotation of the eastern oyster genome  (Gómez-Chiarri et al., 2015) advanced these studies even further, allowing for more in-depth comparisons than were previously possible.
Recent advances in Oxford Nanopore Technologies (ONT) DNA sequencing have allowed for sequencing of ultra-long DNA (ULDNA) from various organisms (Frei et al., 2021; Jain et al., 2018; Lu et al., 2024), but methodology has not yet been optimized for mollusks. Many protocols exist to maximize high molecular weight DNA yields for long-read sequencing (Cahyani et al., 2024; Denton et al., 2024; Gong et al., 2019; Jackson & Howard, 2023; Nishii et al., 2023; M. Strickland et al., 2023; Todorovic et al., 2023), but to date there are no reports of methods for mollusk-specific ultra-long reads nor are there molluscan DNA genome assemblies from ultra-long mollusk DNA. In our experience, the lack of high-quality high molecular weight DNA from mollusks is in great part due to co-extraction of mucopolysaccharides during DNA extraction from mollusk tissues. Mucopolysaccharide-rich tissue promotes the oyster’s ability to successfully filter feed (Beninger & St-Jean, 1997), but is detrimental to extraction of ultra-long, high quality DNA (Adema, 2021). Traditional DNA extraction protocols to remove polysaccharides involve extensive washing and mechanical processing that causes shearing of DNA resulting in short fragments, which eliminates the ability for utilizing ONT’s ULDNA technology. Therefore, optimization of methodology for ULDNA extraction from oysters is needed, warranting further methodological testing.
This study aimed to (1) optimize the extraction and sequencing of ULDNA from eastern oysters and (2) evaluate the potential of eastern oysters as “natural eDNA samplers.” Methodology to successfully isolate high-quality ULDNA from eastern oysters was generated, which required successful sequestering of mucopolysaccharides from oyster DNA without shearing the DNA. Optimization of methods to collect ULDNA from eastern oysters required careful consideration of tissue and modification of washing steps implemented by Brown (1989) (Brown, 1989) to reduce the possibility of DNA release prior to cell lysis. The oyster tissue chosen for preparation of ULDNA was haemocytes, the primary circulating immune cells in mollusks (Anderson et al., 1992; Estrada et al., 2021). Haemocytes are documented to be the most relevant oyster tissue to target for response to parasitic infection (Anderson et al., 1995; Beaven & Paynter, 1999), they are known to express a wide variety of genes and proteins, and their transcriptional profiles have been characterized in a related oyster, Magallana gigas (previously Crassostrea gigas (Salvi & Mariottini, 2017)) (Estrada et al., 2021; Gueguen et al., 2009; Jin et al., 2021). Furthermore, these nucleated cells can be harvested from haemolymph withdrawn with relative ease from the adductor muscle sinus, without sacrificing the organism (Anderson et al., 1992). Lastly, haemocytes are easily isolated single cells that facilitate washing and are far easier to lyse than other tissues commonly used for DNA such as mantle, gill, and adductor muscle.  The washing steps previously documented to reduce mucopolysaccharides in molluscs (Brown, 1989) were optimized for tissue homogenates and involved density gradients. This work aimed to eliminate all homogenization steps and optimize washing. Assembly of the ultra-long reads to both reference genomes was completed to assess success, and unmapped reads were classified to the species level where possible. Results were compared with known species to assess the potential use of oysters as eDNA samplers.

\bibliographystyle{sn-basic}
\bibliography{oyster-uldna}

\printbibliography

\end{document}
